\documentclass[11pt,reqno]{amsart}

\usepackage[T1]{fontenc}
\usepackage{lmodern}
\usepackage{microtype}
\usepackage{amsmath,amsthm,mathtools}
\usepackage{mathrsfs}
\usepackage{enumitem}
\usepackage{hyperref}
\usepackage[nameinlink,noabbrev]{cleveref}
\allowdisplaybreaks[2]
\emergencystretch=2em

\hypersetup{
  hidelinks,
  pdfauthor={Shun Hu},
  pdftitle={Exact Denominator Asymptotics, Phase Geometry, and Mobius Disjointness for a Class of Two-Step Triangular Skew Products}
}

\newtheorem{theorem}{Theorem}[section]
\newtheorem{proposition}[theorem]{Proposition}
\newtheorem{lemma}[theorem]{Lemma}
\newtheorem{corollary}[theorem]{Corollary}
\theoremstyle{definition}
\newtheorem{remark}[theorem]{Remark}

\newcommand{\T}{\mathbb{T}}
\newcommand{\R}{\mathbb{R}}
\newcommand{\C}{\mathbb{C}}
\newcommand{\Z}{\mathbb{Z}}
\newcommand{\PP}{\mathbb{P}}
\newcommand{\e}[1]{e^{2\pi i #1}}
\newcommand{\norm}[1]{\left\lVert #1 \right\rVert}
\newcommand{\abs}[1]{\left\lvert #1 \right\rvert}
\newcommand{\set}[1]{\left\{ #1 \right\}}
\newcommand{\qj}{q_j}
\newcommand{\qjp}{q_{j+1}}
\newcommand{\mobius}{\mu}

\title[Exact Denominator Asymptotics and Phase Geometry]{Exact Denominator Asymptotics, Phase Geometry, and M\"obius Disjointness for a Class of Two-Step Triangular Skew Products}
\author[Shun Hu]{Shun Hu}
\address{The Chinese University of Hong Kong, Shenzhen}
\email{shunhu@link.cuhk.edu.cn}
\date{}
\subjclass[2020]{37B05, 11N64, 37A35}
\keywords{Sarnak conjecture, M\"obius disjointness, rigidity, skew products, Bessel functions}

\begin{document}

\begin{abstract}
We study the original Ces\`aro form of Sarnak's conjecture for the two-step triangular skew
product
\[
\begin{aligned}
T_h(x,y,z)
&=(x+\alpha,\ y+\beta\sin(2\pi x),\ z+h(x,y))
\end{aligned}
\]
on $\T^3$, where $\alpha$ is irrational and $h$ is a real-valued finite-band trigonometric
polynomial. For every irrational $\alpha$ we obtain an explicit denominator-scale asymptotic
formula for the cocycle blocks along the continued-fraction denominators $\qj$ of $\alpha$,
governed by a finite Bessel obstruction vector:
\[
\begin{aligned}
S_{\qj}h(x,y)
&=
\qj \sum_n \Gamma_n(h;z_0)\,\e{ny}\exp\!\bigl(i n z_0 \cos \theta_x\bigr)\\
&\qquad +E_j(x,y).
\end{aligned}
\]
Here $z_0=\pi\beta/\sin(\pi\{\alpha\})$, and the coefficients $\Gamma_n(h;z_0)$ are given by an
explicit finite Bessel expression. If $\alpha$ is Diophantine, then
\[
E_j(x,y)=O(\qj^{-1})
\]
uniformly, and on the resonant variety $\Gamma_n(h;z_0)=0$ this yields polynomial-rate rigidity
and hence the original Ces\`aro M\"obius disjointness by the criterion of
Kanigowski--Lema\'nczyk--Radziwi\l\l. The same denominator law also produces a finite-dimensional
obstruction quotient with an explicit beta-uniform kernel. It further yields a sectorial
rigidity picture for $J_0$-anchored support-two and support-three profiles. For the pure
vertical family
\[
z+A\cos(2\pi n_0 y+\varphi),
\]
the obstruction reduces to the single scalar equation $J_0(n_0 z_0)=0$, yielding an explicit
infinite resonant family. Finally, the normalized Lommel coefficients form a parity-adapted
central-factorial basis in the squared order. This yields exact bounded-bandwidth
reconstruction and a Pad\'e reduction of the normalized phase quotient. On the reduced side,
the remaining ambiguity is organized by multiplier-incidence strata, while on the realized side
the corresponding self-normalized quadratic phase map has an explicitly classified zero-phase
fiber. The first bounded windows are governed by a universal top-order quadratic phase relation,
whose first shadows are an exact conic ambiguity locus and a degree-$1$ saturation phenomenon.
\end{abstract}

\maketitle

\section{Introduction}

\subsection{Setup and mechanism}

We study the two-step triangular skew product
\[
\begin{aligned}
T_h(x,y,z)
&=(x+\alpha,\ y+\beta\sin(2\pi x),\ z+h(x,y))
\end{aligned}
\]
on $\T^3$, where $\alpha$ is irrational and $h$ is a real-valued finite-band trigonometric
polynomial. The natural denominator blocks are
\[
\begin{aligned}
S_q h(x,y)
&:=\sum_{k=0}^{q-1} h(x+k\alpha,\ y+S_k\phi(x)),\\
\phi(x)&:=\beta\sin(2\pi x),
\end{aligned}
\]
so the top cocycle interacts with a nonlinear intermediate layer. The special choice
\[
\phi(x)=\beta\sin(2\pi x)
\]
turns the denominator-time phase into an exact cosine difference. After a Jacobi--Anger
expansion, the whole block becomes an explicit finite Bessel expression. This explicit
denominator law is the basic mechanism of the paper.

Write
\[
\begin{aligned}
h(x,y)
&=\sum_{\abs{m}\le M,\ \abs{n}\le N}\widehat h(m,n)\,\e{mx+ny},\\
a&:=\{\alpha\},\qquad
z_0:=\frac{\pi\beta}{\sin(\pi a)},\qquad
\theta_x:=2\pi x-\pi a.
\end{aligned}
\]
For each vertical frequency $n$ define
\[
\Gamma_n(h;z_0)
:=
\sum_{\abs{m}\le M}\widehat h(m,n)(-i)^m e^{i\pi m a} J_m(n z_0).
\]

\subsection{Results and scope}

The first output is dynamical. For every irrational $\alpha$, \Cref{thm:obstruction} gives the
exact denominator asymptotic. If $\alpha$ is Diophantine, \Cref{thm:quantitative-obstruction}
upgrades it to a uniform $O(\qj^{-1})$ remainder bound, and on the resonant variety
\[
\Gamma_n(h;z_0)=0
\]
\Cref{thm:mobius-resonant} yields the original Ces\`aro M\"obius disjointness.

The second output is geometric. \Cref{thm:obstruction-quotient} identifies a finite-dimensional
obstruction quotient and an explicit beta-uniform kernel. On this quotient,
\Cref{thm:shared-profile-finite} gives shared-profile finiteness,
\Cref{thm:anchor-support-two} closes the $J_0$-anchored support-two sector,
\Cref{thm:anchor-evenodd-second-order-rigidity,thm:anchor-oddodd-second-order-rigidity}
exclude the anchored even-odd and odd-odd support-three sectors at second order, and
\Cref{thm:anchor-eveneven-rigidity} closes the anchored even-even support-three sector at
fifth order. Outside these anchored sectors, \Cref{thm:generic-finite-resonance} shows that the
remaining multi-frequency infinite resonance locus is generically thin.

The third output is bounded-bandwidth phase geometry. \Cref{thm:central-factorial-basis} shows that the
normalized Lommel coefficient polynomials form a parity-adapted central-factorial basis in the
squared order. This yields the bounded-bandwidth reconstruction theorem
\Cref{thm:bounded-bandwidth-reconstruction}. On every bounded even/odd bandwidth window,
\Cref{prop:pade-reconstruction,thm:incidence-geometry,prop:boundary-thinness} reduce the
normalized Pr\"ufer inversion problem to finite-dimensional multiplier-incidence geometry, and
\Cref{prop:zero-phase-fiber} identifies the exact universal zero-phase fiber of the realized
self-normalized quadratic phase map. The first reduced-side windows can already be computed
exactly: \Cref{prop:small-window-incidence} gives a circular degree-$1$ ambiguity locus on
$(1,1)$ and full degree-$1$ saturation from $(1,1)$ into $(2,1)$, both governed by the same
universal top-order quadric of \Cref{prop:top-order-quadric}.

In the pure vertical specialization
\[
\begin{aligned}
T_{A,n_0,\varphi}(x,y,z)
&=(x+\alpha,\ y+\beta\sin(2\pi x),\\
&\qquad z+A\cos(2\pi n_0 y+\varphi)),
\end{aligned}
\]
\Cref{prop:pure-vertical,thm:flagship} reduce the obstruction to the single scalar
$J_0(n_0z_0)$, producing an explicit infinite resonant family.

\subsection{Context and organization}

The closest precedents are the one-step skew-product results of de Faveri, Liu--Sarnak, and Wang
\cite{deFaveri,LiuSarnak,WangAnalytic}, together with the rigidity-to-M\"obius criterion of
Kanigowski, Lema\'nczyk and Radziwi\l\l \cite{KLR}. We use the KLR route, not the
complexity route of Huang, Wang and Ye \cite{HWY}. The genuinely model-specific input is the exact
cosine-difference normal form produced by $\phi(x)=\beta\sin(2\pi x)$. This is also why the
paper does not claim a comparable theorem for a general intermediate cocycle. The Diophantine
assumption enters only at the quantitative stage: the exact denominator law is unconditional,
while the Ces\`aro M\"obius theorem uses Diophantine control of the Dirichlet kernels.

Throughout the paper:
\begin{enumerate}[label=\arabic*.]
\item $\alpha$ is irrational;
\item $\qj$ denotes the continued-fraction denominators of $\alpha$;
\item in the main M\"obius theorems we assume $\alpha$ is Diophantine;
\item $h$ is a real-valued finite-band trigonometric polynomial unless explicitly stated
otherwise;
\item $\widehat h(0,0)=0$;
\item $a=\{\alpha\}$, $z_0=\pi\beta/\sin(\pi a)$, and $\theta_x=2\pi x-\pi a$.
\end{enumerate}

We fix once and for all a compatible translation-invariant metric on $\T$, still denoted by
$d$, and use the associated product metric on $\T^3$. All sup-norm and rigidity statements are
understood with respect to this choice.

The paper follows the proof chain of the denominator law. Section~\ref{sec:normal-form}
derives the denominator-time normal form for the intermediate cocycle. Section~\ref{sec:single-mode}
computes the single-mode denominator blocks by a Jacobi--Anger expansion. Section~\ref{sec:obstruction}
passes from the single-mode calculation to the exact denominator law for finite-band cocycles.
Section~\ref{sec:quantitative} establishes the quantitative $O(\qj^{-1})$ remainder bound under a
Diophantine condition, derives resonant polynomial-rate rigidity, and verifies the precise
Kanigowski--Lema\'nczyk--Radziwi\l\l polynomial-rigidity hypothesis. Section~\ref{sec:mobius}
uses that rigidity input to prove \Cref{thm:mobius-resonant}. Section~\ref{sec:flagship}
treats the pure vertical specialization. Section~\ref{sec:geometry} returns to the obstruction
quotient and develops the support-two and support-three rigidity results together with the
generic finite-resonance theorem. Section~\ref{sec:finite-jet} develops the central-factorial
basis, bounded-bandwidth reconstruction, the reduced-side multiplier-incidence geometry, and the
realized quadratic phase geometry on bounded windows.

\section{Denominator-time normal form}\label{sec:normal-form}

We begin by rewriting the intermediate cocycle at denominator times. The point of this section
is that the special sinusoidal form of $\phi$ turns the denominator-time phase into a cosine
difference, which is the normal form used throughout the remainder of the paper.

\begin{proposition}\label{prop:normal-form}
For every $k\ge 0$ and $x\in\T$,
\[
\begin{aligned}
S_k\phi(x)
&=
\beta\,\Im\!\left(\e{x}\frac{1-\e{k\alpha}}{1-\e{\alpha}}\right)\\
&=
\frac{z_0}{2\pi}\bigl(\cos \theta_x-\cos(\theta_x+2\pi\{k\alpha\})\bigr).
\end{aligned}
\]
Consequently, for every integer $n$,
\[
\begin{aligned}
\e{n S_k\phi(x)}
&=
\exp\!\bigl(i n z_0\cos \theta_x\bigr)\\
&\qquad \times \exp\!\bigl(-i n z_0\cos(\theta_x+2\pi\{k\alpha\})\bigr).
\end{aligned}
\]
\end{proposition}

\begin{proof}
The geometric-series identity gives
\[
\sum_{r=0}^{k-1}\e{x+r\alpha}
=
\e{x}\frac{1-\e{k\alpha}}{1-\e{\alpha}}.
\]
Taking imaginary parts and using $\sin(2\pi u)=\Im(\e{u})$ proves the first identity. Next,
\[
1-\e{t}=-2i\,e^{\pi i t}\sin(\pi t),
\]
so
\[
\frac{1-\e{t}}{1-\e{\alpha}}
=
e^{\pi i(t-a)}\frac{\sin(\pi t)}{\sin(\pi a)}.
\]
Substituting $t=\{k\alpha\}$ yields
\[
S_k\phi(x)
=
\beta \frac{\sin(\pi\{k\alpha\})\sin(\theta_x+\pi\{k\alpha\})}{\sin(\pi a)}.
\]
Using $2\sin u\sin v=\cos(v-u)-\cos(u+v)$ gives the cosine-difference form. Exponentiating gives
the last identity.
\end{proof}

\section{Single-mode denominator blocks}\label{sec:single-mode}

Once the intermediate cocycle is in normal form, the top cocycle can be analyzed one Fourier mode
at a time. The computation in this section isolates the diagonal term responsible for the main
denominator contribution and separates it from the off-diagonal remainder.

\begin{proposition}\label{prop:single-mode}
For the Fourier mode
\[
h_{m,n}(x,y)=\e{mx+ny},
\]
one has
\[
\begin{aligned}
S_q h_{m,n}(x,y)
&=
\e{mx+ny}\exp\!\bigl(i n z_0\cos\theta_x\bigr)\\
&\qquad \times \sum_{r\in\Z} (-i)^r J_r(nz_0)e^{i r\theta_x} D_q((m+r)\alpha),
\end{aligned}
\]
where
\[
D_q(\xi):=\sum_{k=0}^{q-1}\e{k\xi}.
\]
In particular, at $q=\qj$,
\[
\begin{aligned}
S_{\qj} h_{m,n}(x,y)
&=
\qj (-i)^m e^{i\pi m a} J_m(nz_0)\,\e{ny}\exp\!\bigl(i n z_0\cos\theta_x\bigr)\\
&\qquad +R_j^{m,n}(x,y),
\end{aligned}
\]
where
\[
\begin{aligned}
R_j^{m,n}(x,y)
&=
\e{mx+ny}\exp\!\bigl(i n z_0\cos\theta_x\bigr)\\
&\qquad \times \sum_{r\ne -m} (-i)^r J_r(nz_0)e^{i r\theta_x} D_{\qj}((m+r)\alpha).
\end{aligned}
\]
\end{proposition}

\begin{proof}
By definition,
\[
\begin{aligned}
S_q h_{m,n}(x,y)
&=
\sum_{k=0}^{q-1}\e{m(x+k\alpha)+n(y+S_k\phi(x))}.
\end{aligned}
\]
Use \Cref{prop:normal-form} to write
\[
\e{nS_k\phi(x)}
=
\exp\!\bigl(i n z_0\cos\theta_x\bigr)
\exp\!\bigl(-i n z_0\cos(\theta_x+2\pi k\alpha)\bigr).
\]
Apply the Jacobi--Anger expansion
\[
\exp(-iu\cos v)=\sum_{r\in\Z}(-i)^r J_r(u)e^{i r v}
\]
with $u=nz_0$ and $v=\theta_x+2\pi k\alpha$. Interchanging the finite sum in $k$ with the
absolutely convergent Bessel series gives
\[
\begin{aligned}
S_q h_{m,n}(x,y)
&=
\e{mx+ny}\exp\!\bigl(i n z_0\cos\theta_x\bigr)\\
&\qquad \times \sum_{r\in\Z}(-i)^r J_r(nz_0)e^{i r\theta_x} D_q((m+r)\alpha),
\end{aligned}
\]
which is the claimed factorization formula. At $q=\qj$ the unique diagonal index is $r=-m$.
Since $D_{\qj}(0)=\qj$ and $J_{-m}(u)=(-1)^m J_m(u)$, the diagonal contribution is
\[
\begin{aligned}
\qj (-i)^{-m}(-1)^m J_m(nz_0)\,\e{mx+ny} e^{-im\theta_x}
\,\exp\!\bigl(i n z_0\cos\theta_x\bigr).
\end{aligned}
\]
Finally,
\[
\e{mx}e^{-im\theta_x}
=
e^{2\pi i m x}e^{-im(2\pi x-\pi a)}
=
e^{i\pi m a},
\]
so the diagonal term becomes the main term stated in the proposition.
\end{proof}

\section{Exact denominator law}\label{sec:obstruction}

We now pass from the single-mode factorization to the full finite-band cocycle. The key point is
that, at denominator times, each Fourier mode contributes one diagonal term and an explicit
off-diagonal remainder. Summing the diagonal terms produces the obstruction vector.

\begin{theorem}\label{thm:obstruction}
Let
\[
\Gamma_n(h;z_0)
:=
\sum_{\abs{m}\le M}\widehat h(m,n)(-i)^m e^{i\pi m a} J_m(n z_0).
\]
Then
\[
\begin{aligned}
S_{\qj} h(x,y)
&=
\qj \sum_{\abs{n}\le N}\Gamma_n(h;z_0)\,\e{ny}\exp\!\bigl(i n z_0\cos\theta_x\bigr)\\
&\qquad +E_j(x,y),
\end{aligned}
\]
where
\[
E_j(x,y)
:=
\sum_{\abs{m}\le M,\ \abs{n}\le N}\widehat h(m,n) R_j^{m,n}(x,y).
\]
\end{theorem}

\begin{proof}
Expand $h$ into its finitely many Fourier modes and apply \Cref{prop:single-mode} to each mode.
For each pair $(m,n)$, the diagonal contribution arises from the unique index $r=-m$ and
therefore contributes
\[
\begin{aligned}
\qj \widehat h(m,n)(-i)^m e^{i\pi m a} J_m(nz_0)\,\e{ny}
\,\exp\!\bigl(i n z_0\cos\theta_x\bigr).
\end{aligned}
\]
Summing these diagonal terms over the finitely many Fourier modes of $h$ yields the
principal part
\[
\begin{aligned}
\qj \sum_{\abs{n}\le N}\Gamma_n(h;z_0)\,\e{ny}\exp\!\bigl(i n z_0\cos\theta_x\bigr).
\end{aligned}
\]
All terms with $r\ne -m$ are, by definition, absorbed into $E_j$.
\end{proof}

\begin{remark}\label{rem:obstruction-geometry}
The statement of \Cref{thm:obstruction} requires only that $\alpha$ be irrational. No
Diophantine input is used up to this point.
\end{remark}

\begin{remark}\label{rem:obstruction-geometry-2}
Two basic consequences of \Cref{thm:obstruction} will be used repeatedly.
\begin{enumerate}[label=\arabic*.]
\item The zero vertical frequency is the mean:
\[
\Gamma_0(h;z_0)=\widehat h(0,0).
\]
Hence the standing assumption $\widehat h(0,0)=0$ is the vanishing of the
$n=0$ obstruction.
\item If $h$ is real-valued, then
\[
\Gamma_{-n}(h;z_0)=\overline{\Gamma_n(h;z_0)}
\]
for every $n\in\Z$. Thus the obstruction vector is determined by its positive vertical
frequencies.
\end{enumerate}
\end{remark}

\section{Quantitative remainder bounds and resonant rigidity}\label{sec:quantitative}

For the rest of Sections~\ref{sec:quantitative}--\ref{sec:flagship}, assume that $\alpha$ is
Diophantine:
\[
\norm{s\alpha}\ge c_\alpha \abs{s}^{-\tau}
\]
for all nonzero integers $s$.

The goal of this section is to convert the exact obstruction formula of \Cref{thm:obstruction}
into a quantitative denominator-scale rigidity statement. The argument has two ingredients. The
first is a Diophantine bound for the Dirichlet kernels $D_{\qj}(s\alpha)$ on the range
$\abs{s}<\qjp$. The second is the rapid decay of Bessel coefficients in the order, which
controls the complementary tail. Their combination yields an $O(\qj^{-1})$ remainder estimate
for the full finite-band cocycle and therefore polynomial-rate rigidity on the resonant variety.

\begin{lemma}\label{lem:dirichlet}
For every nonzero integer $s$ and every $j$:
\begin{enumerate}[label=\arabic*.]
\item if $0<\abs{s}<\qjp$, then
\[
\abs{D_{\qj}(s\alpha)} \le C_\alpha \abs{s}^{\tau+1}\qjp^{-1};
\]
\item for arbitrary $s$,
\[
\abs{D_{\qj}(s\alpha)} \le \qj.
\]
\end{enumerate}
\end{lemma}

\begin{proof}
Use
\[
D_{\qj}(s\alpha)=\frac{\sin(\pi \qj s\alpha)}{\sin(\pi s\alpha)}.
\]
If $0<\abs{s}<\qjp$, then
\[
\norm{\qj s\alpha}\le \abs{s}\norm{\qj\alpha}\le \abs{s}\qjp^{-1},
\]
so
\[
\abs{\sin(\pi \qj s\alpha)}\le \pi\abs{s}\qjp^{-1}.
\]
The Diophantine condition gives
\[
\abs{\sin(\pi s\alpha)}\ge c'_\alpha \norm{s\alpha}\ge c''_\alpha \abs{s}^{-\tau}.
\]
Combining the estimates proves the first bound. The second follows immediately from the
definition of $D_{\qj}$ as a sum of $\qj$ unimodular terms.
\end{proof}

\begin{lemma}\label{lem:bessel-tail}
Fix $u\in\R$. Then for every integer $r$,
\[
\abs{J_r(u)}\le C_u \frac{(\abs{u}/2)^{\abs{r}}}{\abs{r}!}.
\]
Consequently, for every $A>0$ there exists $C_{u,A}$ such that
\[
\sum_{\abs{r}\ge L} \abs{J_r(u)} \le C_{u,A} L^{-A}
\]
for all integers $L\ge 1$.
\end{lemma}

\begin{proof}
For $r\ge 0$, the series definition of $J_r$ gives
\[
J_r(u)=\sum_{\ell\ge 0} \frac{(-1)^\ell (u/2)^{2\ell+r}}{\ell!(\ell+r)!},
\]
hence
\[
\abs{J_r(u)}
\le
\left(\frac{\abs{u}}{2}\right)^r
\sum_{\ell\ge 0}\frac{(\abs{u}/2)^{2\ell}}{(\ell!)^2}
\le
C_u \frac{(\abs{u}/2)^r}{r!}.
\]
For negative indices use $J_{-r}(u)=(-1)^r J_r(u)$. The tail bound follows because factorial
decay dominates every power.
\end{proof}

\begin{proposition}\label{prop:remainder-one-mode}
For each fixed pair $(m,n)$,
\[
\sup_{x,y}\abs{R_j^{m,n}(x,y)}
\le C_{m,n,\alpha,\beta}\,\qj^{-1}.
\]
\end{proposition}

\begin{proof}
Since the outer phase factors have unit modulus,
\[
\sup_{x,y}\abs{R_j^{m,n}(x,y)}
\le
\sum_{r\ne -m}\abs{J_r(nz_0)}\,\abs{D_{\qj}((m+r)\alpha)}.
\]
Set $s=m+r$ and split into $\abs{s}<\qjp$ and $\abs{s}\ge \qjp$. For the first range,
\Cref{lem:dirichlet} gives
\[
\begin{aligned}
\sum_{0<\abs{s}<\qjp}\abs{J_{s-m}(nz_0)}\,\abs{D_{\qj}(s\alpha)}
&\le
C_\alpha \qjp^{-1}\sum_{s\ne 0}\abs{s}^{\tau+1}\abs{J_{s-m}(nz_0)},
\end{aligned}
\]
and the series on the right converges by \Cref{lem:bessel-tail}. Thus the contribution of this
range is at most $C'_{m,n,\alpha,\beta}\qjp^{-1}$. For the tail use
\[
\qj \sum_{\abs{s}\ge \qjp}\abs{J_{s-m}(nz_0)},
\]
and apply \Cref{lem:bessel-tail} with $A=2$ to obtain
\[
\qj \sum_{\abs{s}\ge \qjp}\abs{J_{s-m}(nz_0)}
\le
C''_{m,n,\beta}\qj \qjp^{-2}
\le
C''_{m,n,\beta}\qj^{-1}.
\]
Combining the two ranges proves the claim.
\end{proof}

\begin{theorem}\label{thm:quantitative-obstruction}
There exists $C(h,\alpha,\beta)$ such that
\[
\sup_{x,y}\abs{E_j(x,y)}
\le C(h,\alpha,\beta)\,\qj^{-1}
\]
for all $j$.
\end{theorem}

\begin{proof}
The sum defining $E_j$ contains only finitely many pairs $(m,n)$, so \Cref{thm:obstruction}
and \Cref{prop:remainder-one-mode} imply
\[
\begin{aligned}
\sup_{x,y}\abs{E_j(x,y)}
&\le
\sum_{\abs{m}\le M,\ \abs{n}\le N}
\abs{\widehat h(m,n)}\,\sup_{x,y}\abs{R_j^{m,n}(x,y)}\\
&\le
C(h,\alpha,\beta)\,\qj^{-1},
\end{aligned}
\]
where the final constant is obtained by summing the finitely many modewise constants.
\end{proof}

\begin{corollary}\label{cor:resonant-cocycle}
If
\[
\Gamma_n(h;z_0)=0
\]
for every active vertical frequency $n$, then
\[
\sup_{x,y}\abs{S_{\qj} h(x,y)} \le C(h,\alpha,\beta)\,\qj^{-1}.
\]
\end{corollary}

\begin{proof}
Under the resonant hypothesis, the principal part in \Cref{thm:obstruction} vanishes
identically, so \Cref{thm:quantitative-obstruction} gives the estimate.
\end{proof}

\begin{corollary}\label{cor:uniform-rigidity}
Under the assumptions of \Cref{cor:resonant-cocycle},
\[
\sup_{(x,y,z)\in \T^3} d(T_h^{\qj}(x,y,z),(x,y,z))
\le C \qj^{-1}.
\]
\end{corollary}

\begin{proof}
The $x$-coordinate displacement is $\norm{\qj\alpha}\le \qj^{-1}$. For the $y$-coordinate,
\Cref{prop:normal-form} gives
\[
\sup_x \abs{S_{\qj}\phi(x)}
\le
\abs{\beta}\,\frac{\abs{1-\e{\qj\alpha}}}{\abs{1-\e{\alpha}}}
\le
C(\alpha,\beta)\,\norm{\qj\alpha}
\le
C(\alpha,\beta)\,\qj^{-1},
\]
while the $z$-coordinate displacement is controlled by \Cref{cor:resonant-cocycle}. Since the
metric on $\T^3$ is the product metric, the three coordinate estimates combine to give the
claimed bound.
\end{proof}

\begin{corollary}\label{cor:klr-condition}
Let $F$ be a trigonometric polynomial on $\T^3$. Then there exists $C_F$ such that
\[
\norm{F\circ T_h^{\qj}-F}_\infty
\le C_F \qj^{-1}.
\]
Moreover, for every $\delta$ with $0<\delta<2/3$,
\[
\sum_{\abs{r}\le \qj^\delta}
\norm{F\circ T_h^{r\qj}-F}_\infty^2 \to 0.
\]
The same conclusion holds in $L^2(\nu)$ for every $T_h$-invariant Borel probability measure
$\nu$.
\end{corollary}

\begin{proof}
Trigonometric polynomials are Lipschitz. By \Cref{cor:uniform-rigidity},
\[
\norm{F\circ T_h^{\qj}-F}_\infty \le C_F \qj^{-1}.
\]
For $\abs{r}\le \qj^\delta$, telescoping gives
\[
\norm{F\circ T_h^{r\qj}-F}_\infty
\le C_F \abs{r}\qj^{-1}\le C_F \qj^{\delta-1}.
\]
Hence
\[
\begin{aligned}
\sum_{\abs{r}\le \qj^\delta}\norm{F\circ T_h^{r\qj}-F}_\infty^2
&\le
(2\qj^\delta+1) C_F^2 \qj^{2\delta-2}\\
&\le
C_F' \qj^{3\delta-2},
\end{aligned}
\]
which tends to $0$ when $\delta<2/3$. The $L^2(\nu)$ statement follows from the sup norm
bound.
\end{proof}

\section{M\"obius disjointness on the resonant variety}\label{sec:mobius}

At this stage the dynamical input is complete: on the resonant variety, the cocycle blocks are
quantitatively rigid along the denominator sequence. This is the hypothesis needed to
invoke the Kanigowski--Lema\'nczyk--Radziwi\l\l criterion.

\begin{theorem}\label{thm:mobius-resonant}
Assume that $\alpha$ is Diophantine and that
\[
\Gamma_n(h;z_0)=0
\]
for every active vertical frequency $n$. Then the system $T_h$ satisfies the original
Ces\`aro Sarnak conjecture.
\end{theorem}

\begin{proof}
The resonant hypothesis implies \Cref{cor:resonant-cocycle}, hence also
\Cref{cor:uniform-rigidity,cor:klr-condition}. Fix a $T_h$-invariant Borel probability measure
$\nu$, and let $\mathcal F_{\mathrm{trig}}$ denote the trigonometric polynomials on $\T^3$.
Then for every $F\in \mathcal F_{\mathrm{trig}}$ and every $\delta\in(0,2/3)$ the denominator
sequence $\qj$ satisfies
\[
\sum_{\abs{r}\le \qj^\delta}
\norm{F\circ T_h^{r\qj}-F}_{L^2(\nu)}^2 \to 0.
\]
Since $\mathcal F_{\mathrm{trig}}$ is linearly dense in $C(\T^3)$, this is exactly the
PR-rigidity condition of \cite[Definition~1.1]{KLR} for the measure-theoretic system
$(\T^3,\mathcal B,\nu,T_h)$. As $\nu$ was arbitrary, the topological system $(\T^3,T_h)$ is
good in the sense of \cite[Definition~1.1]{KLR}. Therefore \cite[Theorem~1.1]{KLR} applies and
yields the M\"obius disjointness of $T_h$, i.e. the original Ces\`aro formulation.
\end{proof}

\section{Pure vertical specialization}\label{sec:flagship}

We conclude by specializing the obstruction formula to a one-frequency family for which the finite
obstruction vector collapses to a single Bessel factor. Because $J_0$ has infinitely many positive
zeros, this produces an infinite explicit resonant family.

For the pure vertical family
\[
h(x,y)=A\cos(2\pi n_0 y+\varphi),
\]
the obstruction vector can be computed explicitly.

\begin{proposition}\label{prop:pure-vertical}
Let
\[
h(x,y)=A\cos(2\pi n_0 y+\varphi).
\]
Then
\[
\Gamma_{n_0}(h;z_0)=\frac{A}{2}e^{i\varphi}J_0(n_0 z_0),
\qquad
\Gamma_{-n_0}(h;z_0)=\frac{A}{2}e^{-i\varphi}J_0(n_0 z_0),
\]
and
\[
\Gamma_n(h;z_0)=0
\]
for every $n\notin\{\pm n_0\}$.
\end{proposition}

\begin{proof}
Write
\[
h(x,y)
=
\frac{A e^{i\varphi}}{2}\,\e{n_0 y}
+
\frac{A e^{-i\varphi}}{2}\,\e{-n_0 y}.
\]
Thus the only nonzero Fourier coefficients of $h$ occur at $(m,n)=(0,n_0)$ and
$(m,n)=(0,-n_0)$. By the definition of $\Gamma_n(h;z_0)$, only the term $m=0$ contributes, and
since $J_0$ is even we obtain
\[
\Gamma_{n_0}(h;z_0)=\frac{A}{2}e^{i\varphi}J_0(n_0 z_0),
\qquad
\Gamma_{-n_0}(h;z_0)=\frac{A}{2}e^{-i\varphi}J_0(n_0 z_0).
\]
All remaining vertical frequencies vanish identically.
\end{proof}

\begin{theorem}\label{thm:flagship}
Let
\[
\begin{aligned}
T_{A,n_0,\varphi}(x,y,z)
&=(x+\alpha,\ y+\beta\sin(2\pi x),\\
&\qquad z+A\cos(2\pi n_0 y+\varphi)).
\end{aligned}
\]
If $\alpha$ is Diophantine and
\[
J_0(n_0 z_0)=0,
\]
then $T_{A,n_0,\varphi}$ satisfies the original Ces\`aro form of Sarnak's conjecture.
\end{theorem}

\begin{proof}
By \Cref{prop:pure-vertical}, the obstruction vector vanishes if and only if
$J_0(n_0 z_0)=0$. The conclusion now follows from \Cref{thm:mobius-resonant}.
\end{proof}

\begin{corollary}\label{cor:infinite-resonant-family}
Fix a Diophantine $\alpha$, an integer $n_0\ge 1$, and parameters $A\neq 0$ and
$\varphi\in\T$. Then there exist infinitely many real values of $\beta$ for which
$T_{A,n_0,\varphi}$ satisfies the original Ces\`aro form of Sarnak's conjecture.
\end{corollary}

\begin{proof}
Let $(\xi_m)_{m\ge 1}$ be the positive zeros of $J_0$. Since $J_0$ has infinitely many positive
zeros, this sequence is infinite. For each $m$, define
\[
\beta_m:=\frac{\xi_m\sin(\pi\{\alpha\})}{\pi n_0}.
\]
Then
\[
n_0 z_0
=
n_0\frac{\pi\beta_m}{\sin(\pi\{\alpha\})}
=
\xi_m,
\]
so $J_0(n_0 z_0)=J_0(\xi_m)=0$. The conclusion follows from \Cref{thm:flagship}.
\end{proof}

\begin{corollary}\label{cor:basic-model}
Consider
\[
\begin{aligned}
T(x,y,z)
&=(x+\alpha,\ y+\beta\sin(2\pi x),\ z+\cos(2\pi y)).
\end{aligned}
\]
If $\alpha$ is Diophantine and
\[
J_0(z_0)=0,
\]
then $T$ satisfies the original Ces\`aro form of Sarnak's conjecture.
\end{corollary}

\begin{proof}
\Cref{thm:flagship} applied with $A=1$, $n_0=1$, and $\varphi=0$ gives the claim.
\end{proof}

\section{Obstruction geometry}\label{sec:geometry}

The exact denominator law of \Cref{thm:obstruction} has a second output beyond the Diophantine
rigidity route: it compresses the finite-band resonance problem to an explicit finite-dimensional
Bessel quotient. The section first records the one-frequency asymptotic input and the resulting
compatibility hierarchy, and then applies them to the quotient. The proofs split into three
types: shared-profile finiteness from compatibility plus formal phase-purity, anchored rigidity
from exact low-order elimination, and generic thinness from placing the compatibility equations on
quotient strata.

\subsection*{The obstruction quotient and its kernel}

Fix bandwidths $(M,N)$. Let $\mathcal H_{M,N}$ be the real vector space of real-valued
trigonometric polynomials
\[
h(x,y)=u_0(x)+\sum_{n=1}^N \bigl(G_n(x)\e{ny}+\overline{G_n(x)\e{ny}}\bigr),
\]
where $u_0$ is a real trigonometric polynomial of degree at most $M$ and
\[
G_n(x)=\sum_{\abs{m}\le M} g_{m,n}\e{mx}.
\]
Define
\[
\mathcal T_{M,N}
:=
\R
\oplus
\bigoplus_{n=1}^N
\operatorname{span}_{\R}
\{J_0(n\cdot),\dots,J_M(n\cdot),\, iJ_0(n\cdot),\dots,iJ_M(n\cdot)\}.
\]
The exact obstruction law defines the real-linear map
\[
\mathcal O_{M,N}:\mathcal H_{M,N}\to \mathcal T_{M,N},
\qquad
h\mapsto \bigl(\Gamma_0(h;\cdot),\Gamma_1(h;\cdot),\dots,\Gamma_N(h;\cdot)\bigr).
\]

\begin{theorem}\label{thm:obstruction-quotient}
The map $\mathcal O_{M,N}$ is surjective. Its kernel is the explicit subspace $\mathcal K_{M,N}$
defined by:
\begin{enumerate}[label=\emph{(\roman*)}]
\item the zero-frequency part $u_0$ has zero mean;
\item for each $1\le n\le N$, the coefficients
\[
d_{m,n}:=g_{m,n}(-i)^m e^{i\pi m a}
\]
satisfy
\[
d_{0,n}=0,
\qquad
d_{-m,n}=-(-1)^m d_{m,n}\quad (1\le m\le M).
\]
\end{enumerate}
Moreover,
\[
\dim_{\R}\mathcal K_{M,N}=2M(N+1),
\qquad
\mathcal H_{M,N}/\mathcal K_{M,N}\cong \mathcal T_{M,N}.
\]
\end{theorem}

\begin{proof}
By \Cref{thm:obstruction,rem:obstruction-geometry-2}, the obstruction splits by vertical
frequency. For $n\ge 1$ one has
\[
\Gamma_n(h;z_0)=\sum_{\abs{m}\le M} d_{m,n} J_m(nz_0).
\]
Using $J_{-m}=(-1)^m J_m$, this vanishes identically if and only if
\[
d_{0,n}=0,
\qquad
d_{-m,n}=-(-1)^m d_{m,n}\quad (1\le m\le M).
\]
For $n=0$, \Cref{rem:obstruction-geometry-2} gives
\[
\Gamma_0(h;z_0)=\widehat u_0(0),
\]
so the zero-frequency obstruction vanishes exactly when $u_0$ has zero mean. This proves the
kernel description.

For the dimension count, the mean-zero $u_0$ sector contributes $2M$ real parameters, and each
positive frequency contributes another $2M$ real parameters, so
\[
\dim_{\R}\mathcal K_{M,N}=2M+N\cdot 2M=2M(N+1).
\]
To prove surjectivity, fix $1\le n\le N$ and let
\[
\Psi_n(t)=\sum_{m=0}^M (a_m+ib_m)J_m(nt)
\]
be an arbitrary element of the $n$th real summand of $\mathcal T_{M,N}$. Define
\[
d_{m,n}:=a_m+ib_m \quad (0\le m\le M),
\qquad
d_{-m,n}:=0 \quad (1\le m\le M),
\]
and then set
\[
g_{m,n}:=d_{m,n} i^m e^{-i\pi m a}
\qquad (\abs{m}\le M).
\]
Let
\[
G_n(x):=\sum_{\abs{m}\le M} g_{m,n}\e{mx},
\qquad
h_n(x,y):=G_n(x)\e{ny}+\overline{G_n(x)\e{ny}}.
\]
Then $h_n\in\mathcal H_{M,N}$ is real-valued by construction. The reality constraint only ties
the $-n$ block to the complex conjugate of the $n$ block; it imposes no relation among the
coefficients $g_{m,n}$ within the positive $n$ block itself. Hence
\[
\Gamma_n(h_n;t)=\sum_{\abs{m}\le M} d_{m,n}J_m(nt)=\Psi_n(t).
\]
So every element of the $n$th summand is realized by a real-valued cocycle. The zero-frequency
summand $\R$ is realized by choosing the mean of $u_0$ arbitrarily. Therefore
$\mathcal O_{M,N}$ is surjective, and the quotient description follows.
\end{proof}

\subsection*{One-frequency asymptotics and compatibility}

The quotient geometry rests on two elementary asymptotic inputs: a one-frequency dichotomy, and
the resulting compatibility relations for any hypothetical common resonant sequence.

\begin{proposition}\label{prop:one-frequency-geometry}
Let
\[
\Gamma(t)=P(t^{-1})J_0(t)+Q(t^{-1})J_1(t),
\qquad
P,Q\in\C[u],
\]
be nontrivial. Write
\[
r:=\min\{\operatorname{ord}_0 P,\operatorname{ord}_0 Q\},
\quad
P(u)=u^r\widetilde P(u),
\quad
Q(u)=u^r\widetilde Q(u),
\]
and set
\[
\alpha:=\widetilde P(0),\qquad \beta:=\widetilde Q(0).
\]
Then:
\begin{enumerate}[label=\emph{(\roman*)}]
\item if
\[
\Im(\overline\alpha\beta)\neq 0,
\]
then $\Gamma$ has only finitely many real zeros;
\item if
\[
\Im(\overline\alpha\beta)=0,
\]
then, after factoring out a common complex phase, the positive zeros of $\Gamma$ admit an
asymptotic expansion
\[
t_k
=
k\pi+\delta+\sum_{r=1}^L \lambda_r(k\pi)^{-r}+O(k^{-L-1})
\]
to every finite order $L$.
\end{enumerate}
\end{proposition}

\begin{proof}
Using the standard Hankel expansions of $J_0$ and $J_1$
\cite[\S10.17,\S10.20]{DLMF}, one has
\[
\Gamma(t)
=
\sqrt{\frac{2}{\pi}}\,t^{-r-\frac12}
\bigl(\alpha\cos\xi+\beta\sin\xi+O(t^{-1})\bigr),
\qquad
\xi:=t-\frac{\pi}{4}.
\]
If $\Im(\overline\alpha\beta)\neq 0$, then the complex oscillatory profile
$\alpha\cos\xi+\beta\sin\xi$ never vanishes on $\R$, hence has a positive lower bound on the
circle. Therefore $\abs{\Gamma(t)}\gg t^{-r-1/2}$ for large $t$, proving~\emph{(i)}.

If $\Im(\overline\alpha\beta)=0$, then $\alpha$ and $\beta$ are real-linearly dependent in
$\C$. Factoring out a common complex phase reduces the leading oscillatory profile to a real
combination $A\cos\xi+B\sin\xi$. The full Hankel series then gives a formal zero equation with
nonvanishing derivative at the leading root. The coefficients of the zero-shift series are then
obtained recursively by coefficient matching, equivalently by applying the implicit function
theorem to the successive truncated equations; compare the standard asymptotic construction of
large Bessel zeros in \cite[\S10.21]{DLMF}. This yields the expansion in~\emph{(ii)}.
\end{proof}

\begin{proposition}\label{prop:compatibility-hierarchy-paper}
Let $\mathcal N$ be a set of active positive frequencies, and for each $n\in\mathcal N$ assume
that $\Gamma_n$ is in the oscillatory-real branch of \Cref{prop:one-frequency-geometry}, with
positive zeros
\[
t_{n,k}
=
k\pi+\delta_n+\sum_{r=1}^L \lambda_{n,r}(k\pi)^{-r}+O(k^{-L-1}).
\]
If there exists an unbounded common real resonant sequence, then for every pair of distinct
active frequencies $m,n\in\mathcal N$ one has
\[
n\delta_m-m\delta_n\in\pi\Z
\]
and
\[
m^{r+1}\lambda_{n,r}=n^{r+1}\lambda_{m,r}
\qquad (1\le r\le L).
\]
\end{proposition}

\begin{proof}
Let $z_j\to\infty$ be the common base parameters, so that $nz_j$ and $mz_j$ lie on the
respective zero sequences. Thus
\[
n z_j
=
k_{n,j}\pi+\delta_n+\sum_{r=1}^L \lambda_{n,r}(k_{n,j}\pi)^{-r}+O(z_j^{-L-1}),
\]
\[
m z_j
=
k_{m,j}\pi+\delta_m+\sum_{r=1}^L \lambda_{m,r}(k_{m,j}\pi)^{-r}+O(z_j^{-L-1}).
\]
Multiply the first identity by $m$ and the second by $n$, then subtract. Since
\[
k_{n,j}=\frac{n}{\pi}z_j+O(1),
\qquad
k_{m,j}=\frac{m}{\pi}z_j+O(1),
\]
comparison of the constant term gives
\[
n\delta_m-m\delta_n\in\pi\Z.
\]
Once this is fixed, comparing the successive powers of $z_j^{-1}$ yields the higher-order
relations inductively.
\end{proof}

\begin{proposition}\label{prop:formal-phase-purity-paper}
Let
\[
F_b(t)=\sum_{m\in S} c_m J_m(t)
\]
be a nontrivial reduced real Bessel profile in the oscillatory-real branch, and let
\[
\mathcal G_b(u):=\sum_{m\in S} c_m (-i)^m H_m(u)
\]
be its formal Hankel amplitude. If all higher zero-shift coefficients of $F_b$ vanish, then
there exists $\phi\in\R$ such that
\[
e^{-i\phi}\mathcal G_b(u)\in\R[[u]].
\]
\end{proposition}

\begin{proof}
Choose the leading phase $\phi$ so that
\[
e^{-i\phi}\mathcal G_b(u)=\mathcal A(u)-i\mathcal B(u),
\]
with $\mathcal A(0)>0$ and $\mathcal B(0)=0$. The formal zero equation near the positive
oscillatory roots is
\[
0=-\mathcal A(u)\sin s+\mathcal B(u)\cos s,
\]
where $s=s(u)\in u\R[[u]]$ is the formal zero-shift series. If all higher zero-shift
coefficients vanish, then $s(u)\equiv 0$, so the equation reduces to $\mathcal B(u)\equiv 0$.
Therefore $e^{-i\phi}\mathcal G_b(u)=\mathcal A(u)$ is real as a formal series.
\end{proof}

\subsection*{Shared-profile multi-frequency finiteness}

We next isolate a large class in which multi-frequency infinite resonance is impossible. This is
the first genuinely global finiteness statement beyond the one-frequency sector.

\begin{theorem}\label{thm:shared-profile-finite}
Let
\[
F_b(t)=\sum_{m\in S} c_m J_m(t)
\]
be a nontrivial reduced real Bessel profile with finite support
$S\subset\Z_{\ge 0}$. Assume that at least two distinct positive active vertical frequencies
carry this same reduced profile, up to nonzero real scalar multiples. Then the common real
resonant parameter set is finite.
\end{theorem}

\begin{proof}
Assume for contradiction that the common real resonant parameter set is infinite. Since the
active frequencies carry the same reduced profile, all of them have the same one-frequency
asymptotic coefficients $\delta,\lambda_1,\lambda_2,\dots$. Applying
\Cref{prop:compatibility-hierarchy-paper} to two distinct active frequencies shows that every
higher coefficient must vanish:
\[
\lambda_r=0
\qquad (r\ge 1).
\]
By \Cref{prop:formal-phase-purity-paper}, there exists $\phi\in\R$ such that
\[
\mathcal G_b(u):=\sum_{m\in S} c_m (-i)^m H_m(u)
\]
has the property that
\[
e^{-i\phi}\mathcal G_b(u)\in \R[[u]].
\]

Write
\[
H_m(u)=E_m(u)+i\,O_m(u),
\]
where $E_m$ is even and $O_m$ is odd in $u$. The imaginary part of
$e^{-i\phi}\mathcal G_b(u)$ therefore yields
\[
\sum_{m\in S} c_m \sin(\phi+m\pi/2)\,E_m(u)=0,
\]
\[
\sum_{m\in S} c_m \cos(\phi+m\pi/2)\,O_m(u)=0.
\]

More concretely, the coefficient of $u^{2r}$ in $E_m$ has the form
\[
\kappa_r m^{4r}+\text{(lower powers of $m^2$)},
\qquad
\kappa_r\neq 0,
\]
while the coefficient of $u^{2r+1}$ in $O_m$ has the form
\[
\kappa_r' m^{4r+2}+\text{(lower powers of $m^2$)},
\qquad
\kappa_r'\neq 0.
\]
Choose distinct indices $m_1,\dots,m_s\in S$. Evaluating the coefficients of
$u^0,u^2,\dots,u^{2s-2}$ in a linear relation among the $E_{m_j}$ gives a square matrix whose
leading part is the Vandermonde matrix
\[
\Bigl((m_j^2)^r\Bigr)_{
\substack{0\le r\le s-1\\ 1\le j\le s}},
\]
which is invertible because the numbers $m_j^2$ are distinct. Hence
$\{E_m\}_{m\in S}$ is linearly independent. The same argument, now using the coefficients of
$u^1,u^3,\dots,u^{2s-1}$, shows that $\{O_m\}_{m\in S}$ is linearly independent as well.

Therefore every active $m\in S$ satisfies
\[
c_m\sin(\phi+m\pi/2)=0,
\qquad
c_m\cos(\phi+m\pi/2)=0,
\]
so $c_m=0$ for every $m\in S$. This contradicts the nontriviality of $F_b$.
\end{proof}

\subsection*{$J_0$-anchored support-two rigidity}

The shared-profile theorem closes the synchronized multi-frequency sector. We now record a
different geometric phenomenon: once one active frequency is the pure vertical anchor $J_0$, the
entire support-two non-shared sector is already closed.

\begin{theorem}\label{thm:anchor-support-two}
Assume that one active positive vertical frequency carries the pure vertical anchor profile
$J_0$. Let $k\ge 1$, and let a second active positive vertical frequency carry a nontrivial
profile supported on $\{0,k\}$:
\[
F(t)=b_0J_0(t)+b_kJ_k(t),
\qquad
(b_0,b_k)\neq (0,0).
\]
If the two active frequencies belong to a common infinite-resonance configuration, then
\[
b_k=0.
\]
Equivalently, the second profile must again be a scalar multiple of $J_0$.
\end{theorem}

\begin{proof}
Write the two active frequencies as $n$ and $p$, and let
\[
q:=\frac{p^2}{n^2}>0.
\]
If $b_0=0$, then the second profile is a pure Bessel mode $J_k$. For even $k$ the first
zero-shift coefficient is negative, while $\lambda_1(J_0)=1/8>0$, so normalized equality at
order~$1$ is impossible. For odd $k$ one has
\[
\lambda_1(J_k)=\frac{4k^2-1}{8},
\]
so normalized equality would force
\[
\frac{p^2}{n^2}=4k^2-1.
\]
But
\[
(2k-1)^2<4k^2-1<(2k)^2,
\]
so $4k^2-1$ is not a square; this is impossible because $p^2/n^2$ is a rational square.

It remains to treat the mixed case $b_0\neq 0$.

If $k$ is odd, write $k=2r+1$,
\[
a:=4k^2-1,
\qquad
y:=(-1)^r\frac{b_k}{b_0}.
\]
A direct Hankel expansion gives
\[
\lambda_1=\frac{1-ay^2}{8(1+y^2)},
\qquad
\lambda_2=
\frac{y(a+1)\bigl((a+9)y^2-(a-7)\bigr)}{128(1+y^2)^2}.
\]
Compatibility with the anchor $J_0$ forces $\lambda_2=0$. Hence either $y=0$, which yields
$b_k=0$, or
\[
y^2=\frac{a-7}{a+9}.
\]
In the latter case normalized equality at order~$1$ gives
\[
q=8\lambda_1=\frac{1-ay^2}{1+y^2}
=-\frac{(a-9)(a-1)}{2(a+1)}<0,
\]
contradiction.

If $k$ is even, write $k=2r$,
\[
a:=4k^2-1,
\qquad
y:=(-1)^r\frac{b_k}{b_0}.
\]
Then the same-parity Hankel expansion yields
\[
\lambda_1=\frac{1-ay}{8(1+y)}.
\]
Compatibility with the anchor implies, after eliminating the third-order coefficient,
\[
y\bigl(A(a)+B(a)y+C(a)y^2\bigr)=0,
\]
where
\[
A(a)=a^3-29a^2-253a-223,
\]
\[
B(a)=-a^3+157a^2-67a-225,
\]
\[
C(a)=192a(a+1).
\]
If $y\neq 0$, then using
\[
q=\frac{1-ay}{1+y}
\]
to eliminate $y$ gives
\[
\begin{aligned}
2(a+1)^3q^2
&+\bigl(3a^4-216a^3-666a^2-672a-225\bigr)q\\
&+a\bigl(a^4-30a^3-96a^2-98a-33\bigr)=0.
\end{aligned}
\]
For $k\ge 4$ one has $a\ge 63$, and all three coefficients above are strictly positive, which is
impossible for $q>0$. The remaining case $k=2$ is an explicit low-order calculation of the same
type and again forces $b_2=0$.

Thus in every case $b_k=0$.
\end{proof}

\subsection*{A second-order normal form in the anchored even-odd support-three sector}

\begin{proposition}\label{prop:anchor-evenodd-second-order}
Let $m\ge 2$ be even and $n\ge 1$ odd. Set
\[
\widetilde J_m:=(-1)^{m/2}J_m,
\qquad
\widetilde J_n:=(-1)^{(n-1)/2}J_n,
\qquad
M:=m^2,
\qquad
N:=n^2.
\]
Consider the reduced profile
\[
F(t)=J_0(t)+b\widetilde J_m(t)+c\widetilde J_n(t).
\]
Then the first two one-frequency shift coefficients of $F$ are
\[
\lambda_1(F)
=
-\frac{M\,b(b+1)+(N-1)c^2}{2\bigl((1+b)^2+c^2\bigr)},
\]
and
\[
\lambda_2(F)
=
-\frac{c\bigl((b+1)\mathcal Q_{M,N}(b)+c^2\mathcal R_{M,N}(b)\bigr)}
{8\bigl((1+b)^2+c^2\bigr)^2},
\]
where
\[
\begin{aligned}
\mathcal Q_{M,N}(b)
&=
N^2-2N+1\\
&\qquad
+b(-M^2-2MN+6M+2N^2-4N+2)\\
&\qquad
+b^2(M^2-2MN+6M+N^2-2N+1),
\end{aligned}
\]
\[
\mathcal R_{M,N}(b)
=
-M^2b+2MNb+2Mb-N^2b-N^2+2Nb+2N-b-1.
\]
\end{proposition}

\begin{proof}
Using the exact Bessel recurrence in the $(J_0,J_1)$ basis, one has
\[
\widetilde J_m(t)
=
\Bigl(1-\frac{M(M-4)}{8}u^2+O(u^4)\Bigr)J_0(t)
\;+\;
\Bigl(-\frac{M}{2}u+O(u^3)\Bigr)J_1(t),
\]
\[
\widetilde J_n(t)
=
\begin{aligned}[t]
&\Bigl(\frac{N-1}{2}u+O(u^3)\Bigr)J_0(t)\\
&\qquad +\Bigl(1-\frac{(N-1)^2}{8}u^2+O(u^4)\Bigr)J_1(t),
\end{aligned}
\qquad
u=t^{-1}.
\]
Hence
\[
F(t)=A(u)J_0(t)+B(u)J_1(t),
\]
with
\[
A(u)
=
(1+b)+\frac{N-1}{2}cu-\frac{M(M-4)}{8}bu^2+O(u^3),
\]
\[
B(u)
=
c-\frac{M}{2}bu-\frac{(N-1)^2}{8}cu^2+O(u^3).
\]
Rotating the leading phase so that the tangent series has vanishing constant term and expanding
that tangent series to order $u^2$ gives the displayed formulas for $\lambda_1(F)$ and
$\lambda_2(F)$.
\end{proof}

\begin{corollary}\label{cor:anchor-evenodd-mobius}
Retain the notation of \Cref{prop:anchor-evenodd-second-order}, and assume
\[
c\neq 0,
\qquad
b\neq -1,
\qquad
\lambda_2(F)=0.
\]
Then:
\begin{enumerate}[label=\emph{(\roman*)}]
\item the genuinely support-three branch is contained in the explicit candidate curve
\[
c^2=-\frac{(b+1)\mathcal Q_{M,N}(b)}{\mathcal R_{M,N}(b)};
\]
\item along that candidate curve, $\lambda_1$ collapses to the fractional-linear function
\[
\lambda_1
=
\frac{-M^2b+4Mb+N^2b+N^2-2Nb-2N+b+1}
{4(Mb-Nb-N+b+1)};
\]
\item its derivative is
\[
\frac{d\lambda_1}{db}
=
\frac{M(N-1)(M-N-3)}
{4(Mb-Nb-N+b+1)^2}.
\]
\end{enumerate}
In particular, for genuine even-odd support-three data one has $M-N-3\neq 0$, so $\lambda_1$
is a strict coordinate on every real branch of the second-order candidate curve.
\end{corollary}

\begin{proof}
Part~\emph{(i)} is the equation $\lambda_2(F)=0$ from
\Cref{prop:anchor-evenodd-second-order}. Substituting that relation into the formula for
$\lambda_1(F)$ yields part~\emph{(ii)}, and differentiating the resulting fractional-linear
expression gives part~\emph{(iii)}. Finally, if $M-N-3=0$, then
\[
(2k)^2-(2r+1)^2=3,
\]
which is impossible because it would factor as
\[
\bigl(2k-(2r+1)\bigr)\bigl(2k+(2r+1)\bigr)=3
\]
with no integer solutions.
\end{proof}

\subsection*{Second-order projective rigidity in the anchored even-odd sector}

\begin{theorem}\label{thm:anchor-evenodd-second-order-rigidity}
Assume that one active positive vertical frequency carries the pure vertical anchor profile
$J_0$. Let a second active positive vertical frequency carry a profile of the form
\[
F(t)=aJ_0(t)+bJ_m(t)+cJ_n(t),
\qquad
a\neq 0,
\qquad
m\ge 2\ \text{even},
\qquad
n\ge 1\ \text{odd}.
\]
If the two active frequencies belong to a common infinite-resonance configuration, then
\[
c=0.
\]
Equivalently, the anchored even-odd support-three sector is disjoint from the $J_0$ fiber of
the normalized Pr\"ufer map already at second order.
\end{theorem}

\begin{proof}
By scaling, we may assume $a=1$ and replace $J_m,J_n$ by the normalized modes
\[
\widetilde J_m:=(-1)^{m/2}J_m,
\qquad
\widetilde J_n:=(-1)^{(n-1)/2}J_n,
\]
which does not affect real solvability. Retain the notation
\[
M:=m^2,
\qquad
N:=n^2.
\]
Suppose first that $c\neq 0$.

If $b=-1$, then \Cref{prop:anchor-evenodd-second-order} gives
\[
\lambda_1(F)=\frac{1-N}{2}<0,
\]
which is incompatible with the anchor normalization
\[
q:=8\lambda_1(F)>0.
\]
So we may also assume $b\neq -1$.

By \Cref{cor:anchor-evenodd-mobius}, the conditions
\[
c\neq 0,
\qquad
b\neq -1,
\qquad
\lambda_2(F)=0
\]
force $F$ to lie on the second-order candidate curve, and along that curve the relation
\[
q=8\lambda_1(F)>0
\]
has the unique solution
\[
b=
\frac{2N^2+Nq-4N-q+2}
{2M^2+Mq-8M-2N^2-Nq+4N+q-2}.
\]
Substituting this into the candidate-curve equation from
\Cref{cor:anchor-evenodd-mobius}\emph{(i)} yields
\[
c^2
=
-\frac{M^2(2M+q-8)\bigl(q^2+(2M+4N-12)q+8(N-1)^2\bigr)}
{(4N+q-4)\bigl(2M^2+Mq-8M-2N^2-Nq+4N+q-2\bigr)^2}.
\]
Now
\[
4N+q-4>0,
\qquad
2M+q-8\ge q>0,
\]
and
\[
q^2+(2M+4N-12)q+8(N-1)^2>0.
\]
Hence the right-hand side is strictly negative, contradicting real-valuedness of $c^2$.

Therefore $c\neq 0$ is impossible, and so $c=0$.
\end{proof}

\subsection*{Second-order projective rigidity in the anchored odd-odd sector}

\begin{theorem}\label{thm:anchor-oddodd-second-order-rigidity}
Assume that one active positive vertical frequency carries the pure vertical anchor profile
$J_0$. Let a second active positive vertical frequency carry a profile of the form
\[
F(t)=aJ_0(t)+bJ_m(t)+cJ_n(t),
\qquad
a\neq 0,
\qquad
m,n\ge 1\ \text{odd},
\qquad
m\neq n.
\]
If the two active frequencies belong to a common infinite-resonance configuration, then
\[
b=c=0.
\]
\end{theorem}

\begin{proof}
By scaling, we may assume $a=1$. Normalize the odd modes by
\[
\widetilde J_m:=(-1)^{(m-1)/2}J_m,
\qquad
\widetilde J_n:=(-1)^{(n-1)/2}J_n,
\]
and write
\[
M:=m^2,
\qquad
N:=n^2.
\]
Set
\[
u:=b+c,
\qquad
d:=(M-1)b+(N-1)c.
\]
In these coordinates the first two normalized Pr\"ufer coefficients are
\[
\lambda_1(F)=-\frac{du}{2(1+u^2)},
\]
\[
\lambda_2(F)=
\frac{2ud^2-(M+N-2)(1+u^2)d+(M-1)(N-1)u(1+u^2)}
{8(1+u^2)^2}.
\]
Now impose the anchor normalization
\[
q:=8\lambda_1(F)>0.
\]
Then
\[
d=-\frac{q(1+u^2)}{4u},
\]
and substituting this into $\lambda_2(F)=0$ gives
\[
u^2
=
-\frac{q\bigl(q+2(M+N-2)\bigr)}
{q^2+2(M+N-2)q+8(M-1)(N-1)}.
\]
The right-hand side is strictly negative for every $q>0$, which is impossible over the reals.
Hence no genuinely support-three anchored odd-odd configuration exists. The remaining
support-two cases are excluded by \Cref{thm:anchor-support-two}, so necessarily $b=c=0$.
\end{proof}

\subsection*{The anchored even-even sector}

\begin{proposition}\label{prop:anchor-eveneven-candidate}
Let
\[
F(t)=J_0(t)+b\widetilde J_m(t)+c\widetilde J_n(t),
\qquad
m,n\ge 2\ \text{even},
\qquad
m\neq n,
\]
where
\[
\widetilde J_\ell:=(-1)^{\ell/2}J_\ell,
\qquad
M:=m^2,
\qquad
N:=n^2.
\]
Set
\[
U:=b+c,
\qquad
D:=Mb+Nc.
\]
Then
\[
\lambda_1(F)=-\frac{D}{2(U+1)},
\qquad
\lambda_2(F)=0.
\]
If moreover
\[
q:=8\lambda_1(F)>0,
\qquad
\lambda_3(F)=-\frac{25}{384}q^2,
\]
then the even-even support-three branch is reduced to the unique candidate curve
\[
U(q)=-q\,\frac{\mathcal N_{M,N}(q)}{\mathcal D_{M,N}(q)},
\]
where
\[
\begin{aligned}[t]
\mathcal N_{M,N}(q)
&=8M^2+8MN+6Mq-64M\\
&\qquad +8N^2+6Nq-64N+q^2-124q+128,
\end{aligned}
\]
\[
\begin{aligned}[t]
\mathcal D_{M,N}(q)
&=32M^2N+8M^2q+32MN^2+32MNq-256MN\\
&\qquad +6Mq^2-64Mq+8N^2q+6Nq^2\\
&\qquad -64Nq+q^3-124q^2+128q,
\end{aligned}
\]
and along that curve one has
\[
\lambda_5(F)-\frac{1073}{5120}q^3
=
\frac{-q\,\mathcal P_{M,N}(q)}
{491520\,(4M+4N+3q-32)},
\]
where $\mathcal P_{M,N}(q)$ is the explicit symmetric quintic
\[
\mathcal P_{M,N}(q)
=
\begin{aligned}[t]
&q^5+(8S-564)q^4+A_{M,N}q^3+B_{M,N}q^2+C_{M,N}q+D_{M,N},
\end{aligned}
\]
with
\[
S:=M+N,
\qquad
P:=MN.
\]
\[
\begin{aligned}
A_{M,N}&=20S^2+20P-2640S+278544,\\
B_{M,N}&=16S^3+88PS-3120S^2-3760P+496000S-3727616,\\
C_{M,N}&=64P^2+96PS^2-1152PS-23552P\\
&\qquad -1984S^3+81920S^2-1068032S+4726784,\\
D_{M,N}&=128P^2S-4096P^2-1024PS^2+36864PS\\
&\qquad -131072P+2048S^3-81920S^2+557056S-1048576.
\end{aligned}
\]
\end{proposition}

\begin{proof}
For even normalized modes one has
\[
\widetilde J_\ell(t)
=
\begin{aligned}
&\Bigl(
1-\frac{\ell^2(\ell^2-4)}{8}u^2
+\frac{\ell^2(\ell^2-4)^2(\ell^2-16)}{384}u^4
+O(u^6)
\Bigr)J_0(t)
\\
&\quad+
\Bigl(
-\frac{\ell^2}{2}u
+\frac{\ell^2(\ell^2-4)^2}{48}u^3
\\
&\qquad\qquad
-\frac{\ell^2(\ell^2-4)^2(\ell^2-16)^2}{3840}u^5
+O(u^7)
\Bigr)J_1(t),
\end{aligned}
\]
Here $u=t^{-1}$.
Substituting these expansions into
\[
F(t)=A(u)J_0(t)+B(u)J_1(t)
\]
and then into the general odd-order Pr\"ufer formulas yields
\[
\lambda_1(F)=-\frac{D}{2(U+1)},
\qquad
\lambda_2(F)=0.
\]
Imposing
\[
q=8\lambda_1(F)>0
\]
gives
\[
D=-\frac{q(U+1)}4.
\]
Substituting this into the order-$3$ relation
\[
\lambda_3(F)=-\frac{25}{384}q^2
\]
produces a linear equation in $U$, whose unique solution is the displayed rational function
$U(q)$. Finally, substituting both
\[
D=-\frac{q(U+1)}4,
\qquad
U=U(q)
\]
into the order-$5$ coefficient gives the displayed residual formula with the explicit symmetric
quintic $\mathcal P_{M,N}(q)$.
\end{proof}

\begin{theorem}\label{thm:anchor-eveneven-rigidity}
Assume that one active positive vertical frequency carries the pure vertical anchor profile
$J_0$. Let a second active positive vertical frequency carry a profile of the form
\[
F(t)=aJ_0(t)+bJ_m(t)+cJ_n(t),
\qquad
a\neq 0,
\qquad
m,n\ge 2\ \text{even},
\qquad
m\neq n.
\]
If the two active frequencies belong to a common infinite-resonance configuration, then
\[
b=c=0.
\]
\end{theorem}

\begin{proof}
By scaling, we may assume $a=1$, and we replace $J_m,J_n$ by the normalized even modes
\[
\widetilde J_m:=(-1)^{m/2}J_m,
\qquad
\widetilde J_n:=(-1)^{n/2}J_n.
\]
Suppose first that
\[
bc\neq 0.
\]
Then we are in the genuinely support-three even-even sector, so
\Cref{prop:anchor-eveneven-candidate} applies. Write
\[
M=4x^2,
\qquad
N=4y^2,
\qquad
2\le x<y,
\]
and let $\mathcal P_{x,y}(q):=\mathcal P_{M,N}(q)$.

We claim that
\[
\mathcal P_{x,y}(q)>0
\qquad (q>0).
\]

If $x\ge 3$, then
\[
\mathcal P_{x,y}(0)=131072(x-1)^2(x+1)^2(y-1)^2(y+1)^2(x^2+y^2-8)>0,
\]
so it is enough to prove
\[
\mathcal P'_{x,y}(q)>0
\qquad (q\ge 0).
\]
Write
\[
a:=x^2,
\qquad
d:=y^2-x^2>0.
\]
Then
\[
\mathcal P'_{x,y}(q)
=
\begin{aligned}[t]
&5q^4+(128x^2+128y^2-2128)q^3\\
&\qquad +48Q_2(a,d)q^2+512Q_1(a,d)q+4096Q_0(a,d),
\end{aligned}
\]
where
\[
Q_2(a,d)=20a^2+20ad-264a+4d^2-132d+17409,
\]
\[
\begin{aligned}
Q_1(a,d)
&=76a^3+114a^2d-1015a^2+46ad^2-1015ad+15500a\\
&\qquad +4d^3-195d^2+7750d-14561,
\end{aligned}
\]
\[
\begin{aligned}
Q_0(a,d)
&=28a^4+56a^3d-284a^3+34a^2d^2-426a^2d+1188a^2\\
&\qquad +6ad^3-204ad^2+1188ad-2086a\\
&\qquad -31d^3+320d^2-1043d+1154.
\end{aligned}
\]
Since $d\ge 2x+1$ and
\[
\partial_dQ_2(a,d)=4(5a+2d-33)>0,
\]
\[
\partial_dQ_1(a,d)=114a^2+92ad-1015a+12d^2-390d+7750>0,
\]
\[
\begin{aligned}[t]
\partial_dQ_0(a,d)
&=56a^3+68a^2d-426a^2+18ad^2-408ad+1188a\\
&\qquad -93d^2+640d-1043>0.
\end{aligned}
\]
throughout the region $x\ge 3$, $d\ge 2x+1$, the coefficients are minimized by the adjacent pair
$y=x+1$. It is therefore enough to check
\[
\Pi_x(q):=\mathcal P'_{x,x+1}(q)>0.
\]
For every $x\ge 3$,
\[
\Pi_x(q)
=
\begin{aligned}[t]
&5q^4+\alpha_x q^3+\beta_x q^2+\gamma_x q+\delta_x,
\end{aligned}
\]
\[
\begin{aligned}
\alpha_x&=256x^2+256x-2128,\\
\beta_x&=48(100x^4+200x^3-1140x^2-1240x+16769),\\
\gamma_x&=512(76x^6+228x^5-717x^4-1814x^3+13799x^2+14744x-7002),\\
\delta_x&=16384(x-1)^2(x+2)^2(7x^4+14x^3-30x^2-37x+25).
\end{aligned}
\]
and every coefficient is strictly positive. Hence $\Pi_x(q)>0$ for $q\ge 0$, so
\[
\mathcal P_{x,y}(q)>0
\qquad (q>0)
\]
throughout the large-even interior.

If $x=2$ and $y\ge 5$, write
\[
r:=y^2-9\ge 16,
\qquad
q=(r+8)z.
\]
After dividing by $(r+8)^2$, the residual becomes
\[
\begin{aligned}
G_r(z)
&=r^3z^5+32r^3z^4+320r^3z^3+1024r^3z^2\\
&\qquad +24r^2z^5+364r^2z^4+1600r^2z^3\\
&\qquad\qquad +12544r^2z^2-28672r^2z\\
&\qquad +192rz^5-320rz^4+199184rz^3\\
&\qquad\qquad +1460224rz^2-4096rz+1179648r\\
&\qquad +512z^5-9472z^4+1654912z^3\\
&\qquad\qquad +16347648z^2+1802240z+5898240.
\end{aligned}
\]
For $r\ge 16$ the coefficients of $z^5$, $z^4$, and $z^3$ are strictly positive, and one can
decompose
\[
G_r(z)=z^3(a_5z^2+a_4z+a_3)+(a_2z^2-bz+c),
\]
with all of $a_5,a_4,a_3$ positive, and
\[
a_2z^2-bz+c>0
\qquad (z>0),
\]
because
\[
a_2=1024r^3+12544r^2+1460224r+16347648,
\]
\[
b=28672r^2+4096r-1802240,
\qquad
c=1179648r+5898240,
\]
\[
4a_2c-b^2
=16777216(239r^4+4954r^3+434487r^2+6652096r+22795280)>0.
\]
Hence $G_r(z)>0$ for all $z>0$, so $\mathcal P_{16,N}(q)>0$ for every $q>0$ whenever
$N\ge 100$.

It remains to treat the isolated pairs $(x,y)=(2,3)$ and $(2,4)$, i.e.
\[
(M,N)=(16,36)
\qquad\text{and}\qquad
(M,N)=(16,64).
\]
Here
\[
\mathcal P_{16,36}(q)=q^5-148q^4+206864q^3+16347648q^2+14417920q+377487360,
\]
\[
\mathcal P_{16,64}(q)=q^5+76q^4+215824q^3+27535104q^2+5529600q+3185049600.
\]
The second polynomial is manifestly positive on $(0,\infty)$. For the first one, a direct
Sturm count shows that it has no positive real root. Since its constant term is positive, it
also satisfies
\[
\mathcal P_{16,36}(q)>0
\qquad (q>0).
\]

Thus $\mathcal P_{M,N}(q)>0$ for every pair of distinct even squares $M,N$ and every $q>0$.
This contradicts the fifth-order anchored compatibility equation from
\Cref{prop:anchor-eveneven-candidate}. Therefore the case $bc\neq 0$ is impossible.

The remaining support-two cases are excluded by \Cref{thm:anchor-support-two}. Hence
\[
b=c=0.
\]
\end{proof}

\subsection*{Generic finite resonance on the quotient}

We finally return to the general multi-frequency quotient geometry. Let $\mathcal N$ be a set of
active positive frequencies with $\abs{\mathcal N}\ge 2$. On a fixed quotient stratum, each
active frequency has a one-frequency reduction
\[
\Gamma_n(t)=P_n(t^{-1})J_0(nt)+Q_n(t^{-1})J_1(nt),
\]
with complex polynomials $P_n,Q_n$, and therefore a leading oscillatory pair
$(\alpha_n,\beta_n)\in\C^2\setminus\{(0,0)\}$. By \Cref{prop:one-frequency-geometry}, the
necessary condition for infinitely many real zeros is
\[
R_n:=\Im(\overline{\alpha_n}\beta_n)=0.
\]
On the corresponding oscillatory-real chart, the positive zeros admit asymptotic expansions
\[
t_{n,k}
=
k\pi+\delta_n+\sum_{r=1}^L \lambda_{n,r}(k\pi)^{-r}+O(k^{-L-1}),
\]
and \Cref{prop:compatibility-hierarchy-paper} forces all normalized coefficients to agree.

\begin{theorem}\label{thm:generic-finite-resonance}
Fix a quotient stratum with at least two active positive frequencies. Then the subset of points
admitting infinitely many common real resonant parameters is contained in a countable union of
proper real-analytic subsets. In particular, on each such stratum the infinite-resonance locus
is meagre and has Lebesgue measure zero.
\end{theorem}

\begin{proof}
If an active frequency $n$ satisfies $R_n\neq 0$, then \Cref{prop:one-frequency-geometry} shows that
$\Gamma_n$ has only finitely many real zeros. Hence infinite common resonance is possible only on
the oscillatory-real chart defined by
\[
R_n=0
\qquad (n\in\mathcal N).
\]
Fix a base frequency $n_1\in\mathcal N$. For each integer vector
\[
\mathbf q=(q_n)_{n\in\mathcal N\setminus\{n_1\}}\in\Z^{\abs{\mathcal N}-1}
\]
and each order $L\ge 1$, let $\mathscr B_{L,\mathbf q}$ be the real-analytic subset cut out by
\[
n\delta_{n_1}-n_1\delta_n=q_n\pi
\qquad (n\neq n_1),
\]
and
\[
n_1^{r+1}\lambda_{n,r}=n^{r+1}\lambda_{n_1,r}
\qquad (n\neq n_1,\ 1\le r\le L).
\]
By \Cref{prop:compatibility-hierarchy-paper}, every point with infinitely many common real resonant
parameters lies in
\[
\bigcup_{\mathbf q}\mathscr B_{L,\mathbf q}
\]
for every $L$.

It remains to show properness. This already follows for $L=1$. On a fixed quotient stratum, each
function $R_n$ depends only on the corresponding one-frequency factor, and the projection onto
that factor is surjective. On that factor there are explicit obstruction data with $R_n\neq 0$,
for example the profile $J_0+iJ_1$. Hence $R_n$ is not identically zero on the stratum, so the
oscillatory-real equation $R_n=0$ cuts out a proper real-analytic subset. Therefore each
$\mathscr B_{1,\mathbf q}$ is proper, and the infinite-resonance locus is contained in a
countable union of proper real-analytic subsets.

The meagreness and measure-zero statements are standard consequences, since a proper
real-analytic subset of a finite-dimensional real-analytic manifold has empty interior and
Lebesgue measure zero.
\end{proof}

\begin{corollary}\label{cor:geometry-landscape}
Let $\Xi\in\mathcal T_{M,N}$ be a quotient datum.
\begin{enumerate}[label=\emph{(\roman*)}]
\item If all positive-frequency obstruction components of $\Xi$ vanish identically, then $\Xi$
lies in the beta-uniform obstruction-free sector.
\item If exactly one positive frequency is active, then the real resonant parameter set is
governed by the one-frequency dichotomy of \Cref{prop:one-frequency-geometry}: it is finite in
the asymptotically nonresonant branch and infinite in the oscillatory-real branch.
\item If at least two distinct positive active frequencies carry the same reduced profile, then
the real resonant parameter set is finite.
\item If one active positive frequency is the anchor $J_0$ and another is supported on
$\{0,k\}$ for some $k\ge 1$, then the real resonant parameter set is finite.
\item If one active positive frequency is the anchor $J_0$ and another lies in the anchored
support-three sector, then the real resonant parameter set is finite. More precisely:
the even-odd and odd-odd sectors are excluded at second order, and the even-even sector is
excluded at fifth order.
\item If at least two positive active frequencies are present in a multi-frequency quotient
stratum, then outside a countable union of proper real-analytic subsets of that stratum the real
resonant parameter set is finite.
\end{enumerate}
\end{corollary}

\begin{proof}
Part~\emph{(i)} is exactly the kernel description in \Cref{thm:obstruction-quotient}.
Part~\emph{(ii)} is \Cref{prop:one-frequency-geometry}. Part~\emph{(iii)} is
\Cref{thm:shared-profile-finite}, part~\emph{(iv)} is \Cref{thm:anchor-support-two},
part~\emph{(v)} is the combination of
\Cref{thm:anchor-evenodd-second-order-rigidity,thm:anchor-oddodd-second-order-rigidity,thm:anchor-eveneven-rigidity},
and part~\emph{(vi)} is \Cref{thm:generic-finite-resonance}.
\end{proof}

\begin{remark}
\Cref{cor:geometry-landscape} describes the current boundary of the problem inside the
finite-dimensional obstruction quotient. The unresolved sector is no longer general
multi-frequency resonance. It is the genuinely non-synchronized, oscillatory-real, effective
multi-frequency sector that survives the explicit kernel, the shared-profile theorem, and the
anchored support-two and support-three rigidity results. In particular, once a $J_0$ anchor is
present, the entire support-three picture is now closed.
\end{remark}

\section{Finite-Jet and Phase Geometry on Bounded Bandwidth Windows}\label{sec:finite-jet}

The support-two and support-three results of \Cref{sec:geometry} are low-order shadows of a
more systematic bounded-bandwidth phase geometry. In this section we isolate the finite-jet and
phase-geometric mechanism behind that picture.

\subsection*{A central-factorial basis for the normalized Lommel coefficients}

For even and odd normalized modes, write
\[
\widetilde J_{2k}(t)=\widetilde A_{2k}(u)J_0(t)+\widetilde B_{2k}(u)J_1(t),
\]
\[
\widetilde J_{2k+1}(t)=\widetilde A_{2k+1}(u)J_0(t)+\widetilde B_{2k+1}(u)J_1(t),
\qquad
u=t^{-1},
\]
with
\[
\widetilde J_{2k}:=(-1)^kJ_{2k},
\qquad
\widetilde J_{2k+1}:=(-1)^kJ_{2k+1}.
\]
These coefficients are governed by the exact Bessel recurrence in the $(J_0,J_1)$ basis, or
equivalently by the classical Lommel polynomials; see \cite[\S10.74]{DLMF}.

\begin{theorem}\label{thm:central-factorial-basis}
For every $k\ge 1$,
\[
\widetilde A_{2k}(u)=\sum_{j=0}^{k-1} E_j\bigl((2k)^2\bigr)u^{2j},
\qquad
\widetilde B_{2k}(u)=\sum_{j=0}^{k-1} F_j\bigl((2k)^2\bigr)u^{2j+1},
\]
and for every $k\ge 0$,
\[
\widetilde A_{2k+1}(u)=\sum_{j=0}^{k} G_j\bigl((2k+1)^2\bigr)u^{2j+1},
\qquad
\widetilde B_{2k+1}(u)=\sum_{j=0}^{k} H_j\bigl((2k+1)^2\bigr)u^{2j}.
\]
The coefficient polynomials are given by
\[
E_0(M)=1,
\]
\[
E_j(M)
=
(-1)^j\,
\frac{
M(M-4)^2(M-16)^2\cdots(M-(2j-2)^2)^2(M-(2j)^2)
}
{4^j(2j)!}
\]
for $j\ge 1$,
\[
F_j(M)
=
(-1)^{j+1}\,
\frac{
M(M-4)^2(M-16)^2\cdots(M-(2j)^2)^2
}
{2\cdot 4^j(2j+1)!},
\]
\[
G_j(N)
=
(-1)^j\,
\frac{
(N-(2j+1)^2)(N-1)^2(N-9)^2\cdots(N-(2j-1)^2)^2
}
{2\cdot 4^j(2j+1)!},
\]
\[
H_0(N)=1,
\qquad
H_j(N)
=
(-1)^j\,
\frac{
(N-1)^2(N-9)^2\cdots(N-(2j-1)^2)^2
}
{4^j(2j)!}
\]
for $j\ge 1$.

In particular, the families $E_j,F_j,G_j,H_j$ form a parity-adapted central-factorial basis in
the squared order: they have strict degree growth and vanish at lower same-parity squares with
the displayed multiplicities.
\end{theorem}

\begin{proof}
The Lommel representation expresses $J_n$ exactly as a combination of $J_0$ and $J_1$ with
polynomial coefficients in $u=t^{-1}$. Substituting $n=2k$ and $n=2k+1$, inserting the sign
normalization, and re-indexing the explicit Lommel sums yields the displayed formulas. The
degree and vanishing statements are immediate from the product structure.
\end{proof}

\subsection*{Bounded-bandwidth reconstruction}

Fix integers $K\ge 1$ and $L\ge 0$, and consider the bounded even/odd window
\[
\mathcal V_{K,L}
:=
\left\{
F(t)=a_0J_0(t)+\sum_{k=1}^{K} b_k \widetilde J_{2k}(t)
+\sum_{\ell=0}^{L} c_\ell \widetilde J_{2\ell+1}(t)
\right\}.
\]
Write
\[
F(t)=A(u)J_0(t)+B(u)J_1(t).
\]
By parity,
\[
B(u)=\sum_{j=0}^{K-1}\beta^{\mathrm{ev}}_j u^{2j+1}
+\sum_{j=0}^{L}\beta^{\mathrm{odd}}_j u^{2j}.
\]
The coefficient of $u^{2K-1}$ in the even block comes only from $\widetilde J_{2K}$, while the
coefficient of $u^{2L}$ in the odd block comes only from $\widetilde J_{2L+1}$. This yields a
global triangular reconstruction.

\begin{theorem}\label{thm:bounded-bandwidth-reconstruction}
For every bandwidth window $(K,L)$, the coefficient map
\[
\Phi_{K,L}:\mathcal V_{K,L}\to \R^{K+L+2},
\]
\[
F\longmapsto
\Bigl(
a_0,\ 
[u^0]B,\ [u^2]B,\dots,[u^{2L}]B,\ 
[u^1]B,\ [u^3]B,\dots,[u^{2K-1}]B
\Bigr)
\]
is a linear isomorphism.

Equivalently, the full coefficient vector
\[
(a_0,b_1,\dots,b_K,c_0,\dots,c_L)
\]
is recovered uniquely from a finite jet of the $J_1$-channel alone.
\end{theorem}

\begin{proof}
By \Cref{thm:central-factorial-basis}, the even contribution to $B(u)$ is triangular in the
support index: the coefficient of $u^{2K-1}$ is exactly $b_KF_{K-1}((2K)^2)$ and this scalar is
nonzero. Hence $b_K$ is recovered first, then recursively $b_{K-1},\dots,b_1$. The same
argument applies to the odd block, where $c_L$ is recovered from the coefficient of $u^{2L}$,
then recursively $c_{L-1},\dots,c_0$. Together with the separated scalar $a_0$, this gives the
inverse to $\Phi_{K,L}$.
\end{proof}

\subsection*{Pad\'e reduction and multiplier-incidence geometry}

For $F\in\mathcal V_{K,L}$, set
\[
z_0:=A(0)+iB(0)\neq 0,
\qquad
C_F(u)+iS_F(u):=\frac{A(u)+iB(u)}{z_0}.
\]
Then
\[
C_F(0)=1,
\qquad
S_F(0)=0,
\qquad
\tan s_F(u)=\frac{S_F(u)}{C_F(u)}=:R_F(u).
\]
Since $A$ and $B$ are polynomials, $R_F$ is a rational function of bounded degree.

\begin{proposition}\label{prop:pade-reconstruction}
Fix $(K,L)$ and define
\[
d_A:=\max(2K-2,\,2L-1),
\qquad
d_B:=\max(2K-1,\,2L).
\]
Then on $\mathcal V_{K,L}$, sufficiently high normalized Pr\"ufer jets determine the rational
phase quotient $R_F=S_F/C_F$ uniquely.
\end{proposition}

\begin{proof}
The formal series $s_F(u)$ determines the jet of $R_F(u)=\tan s_F(u)$ coefficientwise. Now
write
\[
R_F(u)=\frac{P(u)}{Q(u)},
\qquad
Q(0)=1,
\qquad
\deg P\le d_B,\quad \deg Q\le d_A.
\]
The identity $Q(u)R_F(u)-P(u)=0$ modulo $u^{d_A+d_B+2}$ gives the standard square Pad\'e
reconstruction system in the coefficients of $P$ and $Q$, so sufficiently many coefficients of
$R_F$ determine the rational quotient uniquely.
\end{proof}

\begin{theorem}\label{thm:incidence-geometry}
Fix a bounded bandwidth window $(K,L)$ and an integer $s\ge 1$. Let
\[
R(u)=\frac{\widehat S(u)}{\widehat C(u)}
\]
be a normalized rational phase quotient, where $(\widehat C,\widehat S)$ is the unique reduced
coprime representative with $\widehat C(0)=1$ and $\widehat S(0)=0$. Define
\[
\mathcal I_s(K,L)
:=
\left\{
(\widehat C,\widehat S,h):
\deg h=s,\ h(0)=1,\ h(\widehat C,\widehat S)\ \text{is realized in }\mathcal V_{K,L}
\right\},
\]
and let
\[
\mathcal Y_s(K,L)
\]
be its projection to the reduced-pair coordinates.

Then $\mathcal I_s(K,L)$ is a real-algebraic subset of a finite-dimensional affine coefficient
space, and $\mathcal Y_s(K,L)$ is a constructible, hence semialgebraic, subset of the reduced
phase-pair parameter space.
\end{theorem}

\begin{proof}
For fixed $(K,L)$ and $s$, the reduced pair $(\widehat C,\widehat S)$ has bounded degrees
\[
\deg \widehat C\le d_A-s,
\qquad
\deg \widehat S\le d_B-s,
\]
since multiplication by a degree-$s$ polynomial must still fit inside the window bounds
\[
\deg C\le d_A,
\qquad
\deg S\le d_B.
\]
Choose coefficient coordinates for all such reduced pairs, together with coordinates for
\[
h(u)=1+h_1u+\cdots+h_su^s.
\]
The condition that $h(\widehat C,\widehat S)$ be realized in $\mathcal V_{K,L}$ is a finite
system of linear constraints on the coefficients of the product pair, hence a finite polynomial
system in the coefficients of $\widehat C,\widehat S$, and $h$. Coprimeness of
$(\widehat C,\widehat S)$ is the nonvanishing of their resultant, hence a Zariski-open
condition. Therefore $\mathcal I_s(K,L)$ is algebraic on the reduced locus, and its projection
is constructible by elimination.
\end{proof}

\begin{proposition}\label{prop:top-order-quadric}
Let $(C,S)$ be a realized normalized phase pair on a bounded bandwidth window, and let
\[
d:=\deg B
\]
be the top degree of the underlying $J_1$-channel. Write
\[
C(u)=\sum_{j=0}^{d} c_j u^j,
\qquad
S(u)=\sum_{j=1}^{d} s_j u^j.
\]
Then the top two coefficient levels satisfy the universal quadratic relation
\[
2(c_{d-1}s_d-s_{d-1}c_d)+c_d^2+s_d^2=0.
\]
Moreover, this relation is covariant under degree-$1$ multiplication: if
\[
\widetilde C=(1+tu)C,
\qquad
\widetilde S=(1+tu)S,
\]
then
\[
Q_{d+1}(\widetilde C,\widetilde S)=t^2 Q_d(C,S),
\]
where
\[
Q_r(U,V):=2(u_{r-1}v_r-v_{r-1}u_r)+u_r^2+v_r^2
\]
for
\[
U(u)=\sum u_j u^j,
\qquad
V(u)=\sum v_j u^j.
\]
\end{proposition}

\begin{proof}
Write the underlying realized pair as
\[
F(t)=A(u)J_0(t)+B(u)J_1(t),
\qquad
x:=A(0),\quad y:=B(0),\quad D:=x^2+y^2.
\]
If
\[
A(u)=a\,u^{d-1}+\cdots,
\qquad
B(u)=b\,u^d+b' u^{d-1}+\cdots,
\]
then
\[
c_d=\frac{yb}{D},
\qquad
s_d=\frac{xb}{D},
\]
\[
c_{d-1}=\frac{xa+yb'}{D},
\qquad
s_{d-1}=\frac{xb'-ya}{D}.
\]
Therefore
\[
2(c_{d-1}s_d-s_{d-1}c_d)+c_d^2+s_d^2
=
\frac{2ab+b^2}{D}.
\]
In the parity-adapted normalized Lommel basis, the highest coefficients always satisfy
\[
2a+b=0,
\]
so the displayed expression vanishes.

For the covariance statement, if
\[
\widetilde C=(1+tu)C,
\qquad
\widetilde S=(1+tu)S,
\]
then
\[
\widetilde c_{d+1}=t c_d,
\qquad
\widetilde s_{d+1}=t s_d,
\]
\[
\widetilde c_d=c_d+t c_{d-1},
\qquad
\widetilde s_d=s_d+t s_{d-1}.
\]
Substituting these formulas into $Q_{d+1}$ gives
\[
Q_{d+1}(\widetilde C,\widetilde S)=t^2Q_d(C,S).
\]
\end{proof}

\begin{proposition}\label{prop:small-window-incidence}
The first reduced-side degree-$1$ ambiguity loci can be computed exactly.
\begin{enumerate}[label=\emph{(\roman*)}]
\item In the smallest window $(1,0)$, every reduced normalized phase pair has the form
\[
\widehat C(u)=1+pu,
\qquad
\widehat S(u)=qu,
\qquad
p,q\in\R.
\]
Such a pair admits a nontrivial degree-$1$ lift into the window $(1,1)$ if and only if
\[
p^2+(q+1)^2=1.
\]
\item Every reduced normalized pair realized in the window $(1,1)$ admits a nontrivial
degree-$1$ lift into the window $(2,1)$.
\end{enumerate}
\end{proposition}

\begin{proof}
For the window $(1,0)$, write
\[
F=aJ_0+b\widetilde J_2+c\widetilde J_1,
\qquad
\widetilde J_2=J_0-2uJ_1.
\]
Then
\[
A=a+b,
\qquad
B=c-2bu.
\]
With
\[
x:=A(0)=a+b,
\qquad
y:=B(0)=c,
\qquad
D:=x^2+y^2,
\]
the normalized pair is
\[
C(u)=1-\frac{2by}{D}u,
\qquad
S(u)=-\frac{2bx}{D}u.
\]
Hence every pair `(1+pu,qu)` is realized by choosing
\[
x=q,\qquad y=p,\qquad b=-\frac{p^2+q^2}{2},\qquad a=x-b,\qquad c=y.
\]

Now multiply a reduced pair from $(1,0)$ by
\[
h(u)=1+tu,
\qquad
t\neq 0.
\]
The lifted coefficients are
\[
C(u)=1+(p+t)u+pt\,u^2,
\qquad
S(u)=qu+qt\,u^2.
\]
On the window $(1,1)$, \Cref{prop:top-order-quadric} gives exactly the top-order relation
\[
2(c_1s_2-s_1c_2)+c_2^2+s_2^2=0.
\]
Substituting
\[
c_1=p+t,\qquad c_2=pt,\qquad s_1=q,\qquad s_2=qt
\]
gives
\[
t^2\bigl(p^2+(q+1)^2-1\bigr)=0,
\]
which proves part~\emph{(i)}.

For part~\emph{(ii)}, a normalized pair realized in $(1,1)$ has the form
\[
\widehat C(u)=1+c_1u+c_2u^2,
\qquad
\widehat S(u)=s_1u+s_2u^2,
\]
with
\[
2(c_1s_2-s_1c_2)+c_2^2+s_2^2=0.
\]
After multiplication by `1+tu`, the lifted pair becomes
\[
C(u)=1+(c_1+t)u+(c_2+tc_1)u^2+tc_2\,u^3,
\]
\[
S(u)=s_1u+(s_2+ts_1)u^2+ts_2\,u^3.
\]
On the window $(2,1)$, \Cref{prop:top-order-quadric} gives
\[
2(c_2s_3-s_2c_3)+c_3^2+s_3^2=0.
\]
Substituting the lifted coefficients yields
\[
t^2\bigl(2(c_1s_2-s_1c_2)+c_2^2+s_2^2\bigr)=0,
\]
which vanishes identically because the original pair lies in the $(1,1)$ quadric. Hence every
reduced normalized pair from $(1,1)$ admits a degree-$1$ lift into $(2,1)$.
\end{proof}

\begin{proposition}\label{prop:boundary-thinness}
Fix a bounded bandwidth window $\mathcal V_{K,L}$ and let
\[
R(u)=\frac{\widehat S(u)}{\widehat C(u)}
\]
be a normalized rational phase quotient, where $(\widehat C,\widehat S)$ is the unique reduced
coprime representative with $\widehat C(0)=1$. Let
\[
C_F+iS_F=\frac{A_F+iB_F}{z_{0,F}}
\]
be the normalized phase pair of some $F\in\mathcal V_{K,L}$ realizing this quotient, and suppose
\[
C_F=h\,\widehat C,\qquad S_F=h\,\widehat S
\]
for a real polynomial $h$ with $h(0)=1$ and $\deg h=s\ge 0$. Then the top $s$ coefficients of
$S_F$ vanish. More precisely, if
\[
d_B:=\max(2K-1,\,2L),
\]
then
\[
[u^{d_B}]S_F=[u^{d_B-1}]S_F=\cdots=[u^{d_B-s+1}]S_F=0.
\]
\end{proposition}

\begin{proof}
Write
\[
h(u)=1+h_1u+\cdots+h_su^s,
\qquad
h_s\neq 0
\]
if $s\ge 1$, and
\[
\widehat S(u)=\sum_{r\ge 0}\sigma_r u^r.
\]
Since $\deg S_F\le d_B$, the identity
\[
S_F=h\,\widehat S
\]
implies that the coefficients of $u^{d_B+s},u^{d_B+s-1},\dots,u^{d_B+1}$ in $h\,\widehat S$
vanish. The coefficient of $u^{d_B+s}$ is $h_s\sigma_{d_B}$, so $\sigma_{d_B}=0$. Proceeding
downward, the coefficient of $u^{d_B+s-r}$ forces $\sigma_{d_B-r}=0$ for every
$0\le r\le s-1$. Since the top coefficients of $S_F$ agree with those of $\widehat S$ shifted by
the multiplier, this gives the claimed vanishing.
\end{proof}

\subsection*{The realized quadratic phase map}

The realized-side normalized phase geometry is not linear. For
\[
x=(A,B)\in\mathcal V_{K,L},
\]
write
\[
a:=A(0),
\qquad
b:=B(0),
\qquad
z_0:=a+ib.
\]
Then the normalized phase pair is
\[
C_x(u)+iS_x(u):=\frac{A(u)+iB(u)}{a+ib},
\]
so
\[
C_x(u)=\frac{aA(u)+bB(u)}{a^2+b^2},
\qquad
S_x(u)=\frac{aB(u)-bA(u)}{a^2+b^2}.
\]
Accordingly, the realized normalized phase quotient is
\[
R_x(u)=\frac{Q_x(u)}{P_x(u)},
\qquad
P_x(u):=aA(u)+bB(u),
\qquad
Q_x(u):=aB(u)-bA(u).
\]
This is a self-normalized quadratic phase map in the realized coordinates.

\begin{proposition}\label{prop:zero-phase-fiber}
For every bounded bandwidth window $(K,L)$, the zero-phase fiber of the realized quadratic phase
map is exactly
\[
\Psi_{K,L}^{-1}(0)
=
\PP\langle J_0,J_1\rangle
\cup
\PP\langle J_1,\widetilde J_2-J_0\rangle.
\]
\end{proposition}

\begin{proof}
If $R_x\equiv 0$, then $Q_x\equiv 0$. If $A(0)\neq 0$, this forces
\[
B=\lambda A
\]
for some real $\lambda$. The highest-degree triangular terms in the parity-adapted Lommel basis
then eliminate every mode except the constant $J_0$ direction and the constant $J_1$ direction,
so $x\in \PP\langle J_0,J_1\rangle$.

If instead $A(0)=0$, then admissibility implies $B(0)\neq 0$, and $Q_x\equiv 0$ forces
\[
A\equiv 0.
\]
The $A$-channel triangular reconstruction then leaves only the span of $J_1$ and
$\widetilde J_2-J_0$, which gives the second projective line. Both lines are contained in the
fiber by direct substitution.
\end{proof}

\begin{remark}
The bounded-bandwidth inverse problem is therefore no longer a family-by-family elimination
problem. After \Cref{thm:central-factorial-basis,thm:bounded-bandwidth-reconstruction}, the raw
coefficient geometry is globally triangular, and after \Cref{prop:pade-reconstruction}, the
phase quotient is determined by a finite normalized Pr\"ufer jet. On the reduced side, what
remains is multiplier-incidence geometry, and \Cref{prop:boundary-thinness} shows that every
realized nonminimal lift is forced onto a successively thinner top-coefficient boundary stratum.
On the realized side, the first global fiber is already explicit:
\Cref{prop:zero-phase-fiber} shows that the zero-phase fiber is the union of two canonical
projective lines. Thus the bounded-bandwidth frontier has two complementary parts: a
constructible incidence problem for reduced pairs, and a quadratic fiber geometry for realized
pairs.
\end{remark}

\appendix

\section{Verification of the low-order elimination formulas}

The elimination arguments in \Cref{sec:geometry} are all finite and explicit. More precisely,
the anchored support-three theorems use only the low-order coefficients displayed in
\Cref{prop:anchor-evenodd-second-order,prop:anchor-eveneven-candidate} and elementary algebraic
operations on the resulting rational identities. In the mixed-parity and odd-odd sectors, the
second-order contradiction is obtained by direct substitution into the formulas for
$\lambda_1,\lambda_2$. In the even-even sector, the fifth-order obstruction is reduced to the
explicit quintic $\mathcal P_{M,N}(q)$, and the positivity arguments in
\Cref{thm:anchor-eveneven-rigidity} are then proved by the displayed derivative formulas, the
adjacent-pair reduction, and the final Sturm count for the single exceptional polynomial
$\mathcal P_{16,36}$. No computer-assisted search over parameter space is used in the proofs:
the only symbolic input is the exact simplification of the low-order residuals into the
displayed closed formulas.

\begin{thebibliography}{99}

\bibitem{deFaveri}
A.~de Faveri,
\emph{M\"obius disjointness for $C^{1+\varepsilon}$ skew products},
Int. Math. Res. Not. IMRN 2022, no.~4, 2513--2531.
\href{https://doi.org/10.1093/imrn/rnaa185}{doi:10.1093/imrn/rnaa185}.

\bibitem{DLMF}
NIST Digital Library of Mathematical Functions,
F.~W.~J. Olver, A.~B. Olde Daalhuis, D.~W. Lozier, B.~I. Schneider, R.~F. Boisvert,
C.~W. Clark, B.~R. Miller, and B.~V. Saunders (eds.),
\url{https://dlmf.nist.gov/}.

\bibitem{HWY}
W.~Huang, Z.~Wang, and X.~Ye,
\emph{Measure complexity and M\"obius disjointness},
Adv. Math. 347 (2019), 827--858.
\href{https://doi.org/10.1016/j.aim.2019.03.007}{doi:10.1016/j.aim.2019.03.007}.

\bibitem{KLR}
A.~Kanigowski, M.~Lema\'nczyk, and M.~Radziwi\l\l,
\emph{Rigidity in dynamics and M\"obius disjointness},
Fund. Math. 255 (2021), no.~3, 309--336.
\href{https://doi.org/10.4064/fm931-11-2020}{doi:10.4064/fm931-11-2020}.

\bibitem{LiuSarnak}
J.~Liu and P.~Sarnak,
\emph{The M\"obius function and distal flows},
Duke Math. J. 164 (2015), no.~7, 1353--1399.
\href{https://doi.org/10.1215/00127094-2916213}{doi:10.1215/00127094-2916213}.

\bibitem{Sarnak}
P.~Sarnak,
\emph{Three lectures on the M\"obius function randomness and dynamics},
lecture notes, 2011, available at \url{https://publications.ias.edu/sarnak/paper/506}.

\bibitem{WangAnalytic}
Z.~Wang,
\emph{M\"obius disjointness for analytic skew products},
Invent. Math. 209 (2017), no.~1, 175--196.
\href{https://doi.org/10.1007/s00222-016-0707-z}{doi:10.1007/s00222-016-0707-z}.

\end{thebibliography}

\end{document}
